% ============================================================
% P02 — Login Email
% ============================================================
\flowmeta{P02}{Login Email}{Pelanggan}

\begin{flowsection}{Tujuan}
Masuk ke akun menggunakan alamat email/WhatsApp dan kata sandi.
\end{flowsection}

\begin{flowsection}{Prasyarat}
Sudah memiliki akun (lihat alur P01).
\end{flowsection}

\begin{flowsection}{Langkah Login}
\begin{enumerate}
  \item Buka halaman \texttt{/login}.
  \item Isi \textbf{Alamat Email / No. WhatsApp} dan \textbf{Kata Sandi}.
  \item (Opsional) centang \textbf{Ingat saya di perangkat ini}.
  \item Klik \textbf{Masuk Sekarang}.
\end{enumerate}
\end{flowsection}

\shot{gambar/pelanggan/p02-1.png}{Halaman login member TUSKO}{fig:p02-1}

\begin{flowsection}{Hasil yang Diharapkan}
Login berhasil $\rightarrow$ pengguna diarahkan ke beranda dalam keadaan masuk;
navbar menampilkan nama pengguna dan menu akun.
\end{flowsection}

\shot{gambar/pelanggan/p02-2.png}{Pesan error saat email atau kata sandi salah}{fig:p02-2}

\begin{flowsection}{Penanganan Error dan Kasus Khusus}
\begin{itemize}
  \item Kredensial salah $\rightarrow$ ``Email atau kata sandi yang Anda masukkan tidak valid''.
  \item Terlalu banyak percobaan $\rightarrow$ permintaan dibatasi sementara (coba lagi nanti).
  \item Akun nonaktif $\rightarrow$ tidak dapat masuk; hubungi admin.
\end{itemize}
\end{flowsection}

\begin{flowsection}{Kaitan Modul}
FE \texttt{LoginPage.jsx}. BE \texttt{AuthController@login}. Rute
\texttt{POST /api/auth/login}.
\end{flowsection}

\docver{v1.0}{Sep 2026}
